\section{Detailed Validation and Test Metrics}

\subsection{Confusion Matrix and Classification Metrics}

Figure~\ref{fig:confusion_matrix} presents the detailed confusion matrix for the Option B ensemble
on the 56,521-sample held-out test set. The confusion matrix clearly shows the model's classification
performance:

\begin{table}[h!]
\centering
\caption{Confusion Matrix (Test Set, $n=56,521$)}
\label{tab:confusion_matrix}
\begin{tabular}{lrr}
\hline
 & \textbf{Predicted Background} & \textbf{Predicted CWD} \\
\hline
\textbf{Actual Background} & 16,723 (TN) & 294 (FP) \\
\textbf{Actual CWD} & 232 (FN) & 39,272 (TP) \\
\hline
\end{tabular}
\end{table}

\subsubsection{Primary Classification Metrics}

\begin{itemize}
    \item \textbf{Accuracy}: 0.9907 (99.07\%) — Overall proportion of correct predictions
          \[
              \text{Accuracy} = \frac{TP + TN}{TP + TN + FP + FN} = \frac{39,272 + 16,723}{56,521} = 0.9907
          \]

    \item \textbf{Sensitivity (Recall/TPR)}: 0.9941 (99.41\%) — Ability to correctly identify CWD
          \[
              \text{Sensitivity} = \frac{TP}{TP + FN} = \frac{39,272}{39,272 + 232} = 0.9941
          \]
          \emph{Interpretation}: Of 39,504 actual CWD samples, the model correctly identifies 39,272
          (misses only 232). Excellent CWD detection rate.

    \item \textbf{Specificity (TNR)}: 0.9827 (98.27\%) — Ability to correctly identify background
          \[
              \text{Specificity} = \frac{TN}{TN + FP} = \frac{16,723}{16,723 + 294} = 0.9827
          \]
          \emph{Interpretation}: Of 17,017 actual background samples, the model correctly identifies 16,723
          (false alarms on only 294). Excellent background rejection.

    \item \textbf{Precision (PPV)}: 0.9926 (99.26\%) — Confidence when predicting CWD
          \[
              \text{Precision} = \frac{TP}{TP + FP} = \frac{39,272}{39,272 + 294} = 0.9926
          \]
          \emph{Interpretation}: When the model predicts CWD, it is correct 99.26\% of the time.
          Very few false alarms.

    \item \textbf{Negative Predictive Value (NPV)}: 0.9863 (98.63\%) — Confidence when predicting background
          \[
              \text{NPV} = \frac{TN}{TN + FN} = \frac{16,723}{16,723 + 232} = 0.9863
          \]
          \emph{Interpretation}: When the model predicts background, it is correct 98.63\% of the time.
\end{itemize}

\subsubsection{Harmonic and Balanced Metrics}

\begin{itemize}
    \item \textbf{F1 Score}: 0.9933 — Harmonic mean of precision and recall
          \[
              F_1 = 2 \times \frac{\text{Precision} \times \text{Recall}}{\text{Precision} + \text{Recall}}
                  = 2 \times \frac{0.9926 \times 0.9941}{0.9926 + 0.9941} = 0.9933
          \]
          \emph{Interpretation}: Extremely high F1 score indicates excellent balance between
          precision and recall. The model performs equally well on both classes.

    \item \textbf{Matthews Correlation Coefficient (MCC)}: 0.9779 — Balanced measure for binary classification
          \[
              \text{MCC} = \frac{TP \cdot TN - FP \cdot FN}{\sqrt{(TP+FP)(TP+FN)(TN+FP)(TN+FN)}} = 0.9779
          \]
          \emph{Interpretation}: Ranges from -1 to +1. Value of 0.9779 indicates excellent agreement
          between predictions and observations. Not affected by class imbalance.
\end{itemize}

\begin{figure}[H]
\centering
\includegraphics[width=0.95\textwidth]{./output/thesis_visualizations/confusion_matrix_detailed.png}
\caption[Confusion Matrix and Detailed Metrics]{
    Detailed confusion matrix for Option B ensemble on 56,521-sample test set. Left: Heatmap showing
    39,272 true positives, 16,723 true negatives, 294 false positives, 232 false negatives. Right:
    Comprehensive metrics table including TP, TN, FP, FN, and all derived metrics (Accuracy,
    Precision, Recall, Specificity, NPV, F1-Score, Matthews Correlation Coefficient).
}
\label{fig:confusion_matrix}
\end{figure}

\subsection{ROC and Precision-Recall Curves}

Figure~\ref{fig:roc_pr} presents the Receiver Operating Characteristic (ROC) curve and
Precision-Recall (PR) curve, which visualize model performance across all possible decision thresholds.

\begin{figure}[H]
\centering
\includegraphics[width=0.95\textwidth]{./output/thesis_visualizations/roc_pr_curves.png}
\caption[ROC and Precision-Recall Curves]{
    Option B ensemble performance curves. Left: ROC curve with AUC=0.9885, showing excellent
    discrimination between classes across all thresholds. The red dot marks the optimal operating
    point. Right: Precision-Recall curve with AUC=0.8964, demonstrating high precision at all
    recall levels. The gray dashed line shows the positive class baseline (70\%).
}
\label{fig:roc_pr}
\end{figure}

\subsubsection{Interpretation}

\begin{itemize}
    \item \textbf{AUC-ROC (Area Under Curve)}: 0.9885 — Probability that the model ranks a
          random positive example higher than a random negative example. A score of 0.9885 is
          exceptional (scale: 0.5 = random, 1.0 = perfect).

    \item \textbf{AUC-PR}: 0.8964 — Area under the precision-recall curve. Particularly important
          for imbalanced datasets. The ensemble maintains high precision across all recall levels.

    \item \textbf{Optimal threshold}: The red point on the ROC curve marks $t^* = 0.40$, where
          sensitivity minus false positive rate (Youden's J statistic) is maximized.
\end{itemize}

\subsection{Threshold Analysis and Performance Across Decision Boundaries}

Figure~\ref{fig:threshold_analysis} demonstrates how model performance varies across different
classification thresholds. This analysis is critical for understanding trade-offs in the
sensitivity-specificity spectrum.

\begin{figure}[H]
\centering
\includegraphics[width=0.95\textwidth]{./output/thesis_visualizations/threshold_analysis.png}
\caption[Threshold Analysis: Performance Metrics Across Decision Boundaries]{
    Comprehensive analysis of model performance at thresholds from 0 to 1. Top-left: Sensitivity
    (red) and specificity (blue) reveal the trade-off between detecting CWD and rejecting
    background. Top-right: Precision and recall show the inverse relationship (higher threshold
    increases precision but decreases recall). Bottom-left: F1 score and accuracy peak at
    $t^* = 0.40$. Bottom-right: Matthews Correlation Coefficient (MCC) provides a balanced view
    and confirms optimality at $t = 0.40$. The vertical green line marks the optimal threshold.
}
\label{fig:threshold_analysis}
\end{figure}

\subsubsection{Key Observations}

\begin{enumerate}
    \item \textbf{Sensitivity-Specificity Trade-off}: At threshold 0.40, sensitivity (99.41\%) and
          specificity (98.27\%) are well-balanced. Lowering the threshold increases sensitivity
          (more CWD detected) but decreases specificity (more false alarms). Raising the threshold
          has the opposite effect.

    \item \textbf{Precision-Recall Inverse Relationship}: As threshold increases, precision increases
          (fewer false positives among positive predictions) but recall decreases (fewer positives
          detected overall). The threshold 0.40 achieves optimal balance with precision=99.26\% and
          recall=99.41\%.

    \item \textbf{F1 Score Peak}: F1 score is maximized at threshold 0.40, confirming this as the
          optimal operating point for balancing precision and recall.

    \item \textbf{MCC Robustness}: The Matthews Correlation Coefficient, which is robust to class
          imbalance, also peaks at $t = 0.40$, providing independent confirmation of optimality.
\end{enumerate}

\subsection{Calibration Analysis}

Figure~\ref{fig:calibration} shows how well the model's probability estimates match the true
frequency of positive outcomes (calibration curve, also called a reliability diagram).

\begin{figure}[H]
\centering
\includegraphics[width=0.85\textwidth]{./output/thesis_visualizations/calibration_curve.png}
\caption[Calibration Curve (Reliability Diagram)]{
    Calibration curve for Option B ensemble. Each point represents a bin of probabilities. For
    example, when the model predicts probability 0.7, the actual positive rate in that bin should
    be approximately 0.7. The ensemble predictions follow the diagonal (perfect calibration) closely,
    indicating that predicted probabilities accurately reflect true occurrence rates. Slight
    overconfidence at high probabilities (points above the diagonal).
}
\label{fig:calibration}
\end{figure}

\subsubsection{Interpretation}

A perfectly calibrated model would have all points on the diagonal. The Option B ensemble:
\begin{itemize}
    \item Shows excellent overall calibration, with points closely following the diagonal,
    \item Demonstrates slight overconfidence at high predicted probabilities (points above diagonal),
    \item Is well-calibrated for moderate probabilities (0.2–0.7), which is the most decision-critical range.
\end{itemize}

\subsection{Error Analysis: False Positives and True Negatives}

Figure~\ref{fig:anomaly} analyzes the distribution of errors and correct rejections, identifying
which predictions are most uncertain or most confident.

\begin{figure}[H]
\centering
\includegraphics[width=0.95\textwidth]{./output/thesis_visualizations/anomaly_analysis.png}
\caption[False Positive and True Negative Analysis]{
    Error analysis of Option B ensemble. Top-left: Distribution of 294 false positives (background
    predicted as CWD) has mean probability 0.589, showing these are mostly mid-range confidence
    predictions. Top-right: Distribution of 16,723 true negatives (correctly identified background)
    has mean probability 0.127, showing confident rejections. Bottom-left: Top 20 false positives
    ranked by confidence (highest probability), identifying cases where the model was most wrongly
    confident. Bottom-right: Top 20 true negatives ranked by confidence, showing the most confident
    correct background identifications.
}
\label{fig:anomaly}
\end{figure}

\subsubsection{Key Findings}

\begin{itemize}
    \item \textbf{False Positives ($n=294$)}: Mean predicted probability is 0.589, indicating these
          errors typically occur at moderate confidence levels, not extreme overconfidence. This is
          favorable: the model is not highly confident when wrong.

    \item \textbf{True Negatives ($n=16,723$)}: Mean predicted probability is 0.127, showing the
          model typically makes confident background predictions. These are reliable negative
          classifications.

    \item \textbf{Error Rate Analysis}: Only 294 false positives out of 56,521 test samples (0.52\%)
          and only 232 false negatives (0.41\%). Both error rates are extremely low, demonstrating
          excellent generalization to the held-out test set.
\end{itemize}

\subsection{Comprehensive Metrics Summary Table}

\begin{table}[H]
\centering
\caption{Comprehensive Option B Ensemble Metrics (Test Set, $n=56,521$)}
\label{tab:comprehensive_metrics}
\begin{tabular}{llr}
\hline
\textbf{Metric Category} & \textbf{Metric} & \textbf{Value} \\
\hline
\multirow{4}{*}{\textbf{Confusion Matrix}} & True Positives (TP) & 39,272 \\
& True Negatives (TN) & 16,723 \\
& False Positives (FP) & 294 \\
& False Negatives (FN) & 232 \\
\hline
\multirow{5}{*}{\textbf{Classification}} & Accuracy & 0.9907 (99.07\%) \\
& Sensitivity (Recall/TPR) & 0.9941 (99.41\%) \\
& Specificity (TNR) & 0.9827 (98.27\%) \\
& Precision (PPV) & 0.9926 (99.26\%) \\
& NPV & 0.9863 (98.63\%) \\
\hline
\multirow{3}{*}{\textbf{Harmonic \& Balanced}} & F1 Score & 0.9933 \\
& Matthews Correlation Coeff & 0.9779 \\
& AUC-ROC & 0.9885 \\
\hline
\multirow{2}{*}{\textbf{Data Distribution}} & Positive class (CDW) & 39,504 (69.9\%) \\
& Negative class (Background) & 17,017 (30.1\%) \\
\hline
\end{tabular}
\end{table}

\subsection{Class-Specific Performance Analysis}

The model demonstrates excellent performance on both positive (CWD) and negative (background) classes:

\subsubsection{CWD Detection (Positive Class)}

\begin{itemize}
    \item \textbf{Detection Rate}: 99.41\% (39,272 out of 39,504 detected)
    \item \textbf{Misses}: Only 232 CWD samples incorrectly classified as background
    \item \textbf{Confidence}: When model predicts CWD, it is correct 99.26\% of the time
    \item \textbf{Implication}: Excellent for applications where missing CWD is costly (e.g., forest
          inventory, ecological assessment)
\end{itemize}

\subsubsection{Background Rejection (Negative Class)}

\begin{itemize}
    \item \textbf{Rejection Rate}: 98.27\% (16,723 out of 17,017 correctly rejected)
    \item \textbf{False Alarms}: Only 294 background samples incorrectly classified as CWD
    \item \textbf{Confidence}: When model predicts background, it is correct 98.63\% of the time
    \item \textbf{Implication}: Excellent for applications where false positives are costly (e.g.,
          resource allocation, automated flagging systems)
\end{itemize}

\subsection{Summary: What These Metrics Mean}

The Option B ensemble achieves exceptional performance across all metrics:

\begin{enumerate}
    \item \textbf{Accuracy (99.07\%)}: 99 out of 100 predictions are correct. Exceptional overall performance.

    \item \textbf{Sensitivity (99.41\%)}: Detects 99 out of 100 CWD samples. Loss of only 0.59\% of
          target class.

    \item \textbf{Specificity (98.27\%)}: Correctly rejects 98 out of 100 background samples. Only
          1.73\% false alarm rate.

    \item \textbf{Precision (99.26\%)}: When flagging a sample as CWD, the model is correct 99 times out of 100.

    \item \textbf{F1 Score (0.9933)}: Excellent balance between detection (recall) and accuracy (precision).

    \item \textbf{Matthews Correlation (0.9779)}: Near-perfect agreement between predictions and
          true labels, even accounting for class imbalance.
\end{enumerate}

These metrics demonstrate that the Option B ensemble is exceptionally well-suited for practical
deployment in CWD detection applications.
